\chapter{Постановка задачи} \label{ch2}

В данной главе приведена формальная постановка задачи приоритизации тестовых случаев и определены метрики
оценки эффективности. В параграфе~\ref{ch2:sec1} дано формальное определение задачи. В параграфе~\ref{ch2:sec2} 
описаны метрики APFD и APFDc.

\section{Формальное определение} \label{ch2:sec1}

Пусть $T = \{t_1, t_2, \ldots, t_n\}$ --- тестовый набор из $n$ тестовых случаев. 
Каждый тестовый случай $t_i$ характеризуется множеством покрываемых методов $M_i \subseteq M$, вероятностью 
дефектов $r_i \in [0, 1]$ и временем выполнения $c_i > 0$.

Задача приоритизации состоит в нахождении перестановки $T^*$: $$T^* = \arg\max_{T' \in \Pi(T)} f(T')$$
где $\Pi(T)$ --- множество всех перестановок $T$, целевая функция: 
$$f(T') = w_c \cdot \text{Cov}(T') + w_r \cdot \text{FP}(T') + w_d \cdot \text{Div}(T'), \quad w_c + w_r + w_d = 1$$

\section{Метрики оценки эффективности} \label{ch2:sec2}

Основной метрикой оценки является APFD (Average Percentage of Faults 
Detected)~\cite{elbaum2002}:
\begin{equation}\label{eq:apfd}
\text{APFD} = 1 - \frac{\sum_{i=1}^{m} TF_i}{n \cdot m} + \frac{1}{2n}
\end{equation}
где $n$ --- количество тестов, $m$ --- количество дефектов, $TF_i$ --- позиция теста, первым обнаружившего
$i$-й дефект. Значение APFD близкое к~1 означает, что дефекты обнаруживаются в начале выполнения набора.

Метрика APFDc расширяет APFD с учётом стоимости выполнения каждого теста:
$$\text{APFDc} = \frac{\sum_{i=1}^{m} \left( t_{TF_i} - \frac{c_{TF_i}}{2} \right)}{T_{total} \cdot m}$$
где $t_{TF_i}$ --- суммарное время выполнения тестов до момента обнаружения $i$-го дефекта, $c_{TF_i}$ --- 
время выполнения теста обнаружившего дефект, $T_{total}$ --- суммарное время выполнения всего набора.